% META: {"id": "TSPE-INTEG-EX-004", "chapitre": "TSPE-CALCUL-INTEGRAL", "type_objet": "exercice", "capacites": ["TSPE-INTEG-C4"], "capacites_codes": ["C4"], "methodes": ["M2"], "parcours": 2, "competences": ["calculer", "raisonner"], "duree_min": 14, "mode_creation": "ex_nihilo", "sources_inspiration": [], "fichier_tex": "chapitres/TSPE-CALCUL-INTEGRAL/exercices/TSPE-INTEG-EX-004.tex", "corrige_tex": "chapitres/TSPE-CALCUL-INTEGRAL/corriges/TSPE-INTEG-CO-004.tex", "status": "approved"}
% BEGIN-VERIFY
% from sympy import symbols, integrate, Rational
% x = symbols('x')
% assert integrate(x - x**2, (x, 0, 1)) == Rational(1, 6)
% END-VERIFY
\begin{exercice}{TSPE-INTEG-EX-004}{2}{14}
On considere les courbes de $f(x)=x$ et $g(x)=x^2$.
\begin{enumerate}
  \item Determiner les points d'intersection des deux courbes.
  \item Justifier que $f \geq g$ sur $[0,1]$.
  \item Calculer l'aire du domaine compris entre les deux courbes sur $[0,1]$.
\end{enumerate}
\end{exercice}
